\subsection*{从表达式到值编号}

在当前实现中，值编号以语法制导的方式进行：遍历抽象语法树的节点，遇到含表达式的节点就生成值编号。

生成值编号的方法很简单：把表达式作为输入，返回它对应的值编号。

更具体地说，它：

\begin{itemize}[leftmargin=1.6em]
\item 先由表达式生成一个签名。
\item 如果这个签名已经存在，就查询值编号表，返回与它关联的值编号。
\item 否则，在值编号表里把这个签名和一个新值编号关联起来，再返回这个新编号。
\end{itemize}

为简化实现，目前签名就是表达式的字符串；以后可以像多数值编号过程那样改用哈希来改进。

此外，为了处理符号的定义和引用，它遵循以下规则：

\begin{itemize}[leftmargin=1.6em]
\item 表达式用于初始化某个符号（变量或部分类型）时，该符号获得与表达式相同的值编号。更具体地说，它把\emph{带作用域}的符号名签名和值编号关联到表中。
\item 表达式引用某个符号时，就在值编号表里查找对应的签名并返回。
\end{itemize}

注意，生成签名时，符号必须使用\emph{带作用域的名字}，因为不同作用域里允许定义同名符号。
